(1) \(e^{i\phi}=1+i\phi-\phi^2/2+O(\phi^3)\)。平均 \(\langle\phi\rangle=0\)，所以零级参考场近似 \(E_0\)，非零空间频率散射场一阶 \(iE_0\phi(x)\)。
(2) 理想相位板后场
\[
E_i(x)=E_0[ae^{i\delta}+i\phi(x)]+O(\phi^2).
\]
强度一阶
\[
I(x)=|E_0|^2[a^2+2a\phi(x)\sin\delta]+O(\phi^2).
\]
因为 \(2\Re[(ae^{i\delta})^*i\phi]=2a\phi\sin\delta\)，正的相移 \(\delta=\pi/2\) 将正相位变为亮对比。
(3) 以背景强度 \(I_b=a^2|E_0|^2\) 归一，\((I-I_b)/I_b\simeq(2/a)\phi\sin\delta\)。\(|\sin\delta|=1\) 时线性响应最大，\(\delta=\pm\pi/2\) 给正负相衬。若a趋0，背景被切断，虽然相对系数形式发散，实际信噪与线性展开不能保持；非零平均相位应归入参考复振幅并重新定义相位差。
(4) 简化写法 \(|ae^{i\delta}+i\phi-\phi^2/2|^2\) 给二阶项 \((1-a\cos\delta)\phi^2\)（忽略样品二阶零级重新分配时），说明强相位时非线性。严格能量守恒需把二阶散射的傅里叶零级也通过板处理，故局域二阶项并非通用的真实相衬成像定律。吸收和有限孔径引入复振幅及非局域点扩散卷积。
